\documentclass{swufethesis}

\usepackage{zhlipsum}
\usepackage{lipsum}
\usepackage{booktabs}
\usepackage{hyperref}

\addbibresource{sample.bib}

\swufesetup{%
  date = 2021-01-04,
  title = 总之这是一个很长很长很长很长很长很长的——《题目》,
  author = 某某某\textperiodcentered 某某,
  school = 经济信息工程学院,
  discipline = 这也是我财一个很长很长很长很长很长的专业名称（长到要加备注说明）,
  student-id = 4XXXXXXX,
  supervisor = XX老师,
}

\begin{document}
\maketitle % 封面
\statement % 版权申明

\begin{abstract}{毕业论文, 毕业设计, 西南财经大学}
  \zhlipsum[1-4]
\end{abstract}

\begin{abstract*}{Dissertation, Graduation project, Swufe}
  \lipsum[1]
\end{abstract*}

\tableofcontents

\mainmatter % \mainmatter会重置页码样式为阿拉伯数字，并使\chapter被编号
\chapter{背景介绍}
\zhlipsum
\chapter{主体内容}

\section{二级标题}
\subsection{三级标题}
\zhlipsum[2-5]
巴拉两阿里\footnote{注释}，一个英文文献\cite{dirac}提出了，一个中文文献\cite{蔡敏2006--}提出了balabala。

\zhlipsum

\begin{table}[htb]
  \caption{表格示例}
  \centering
  \begin{tabular}{lrrrr}
    \toprule
    项目 & 指标  \\
    \midrule
    A    & 1.7   \\
    B    & 1.2   \\
    C    & -12.6 \\
    \bottomrule
  \end{tabular}
\end{table}

\zhlipsum

\begin{figure}[htb]
  \centering
  \includegraphics[width=0.4\textwidth]{figures/example-image-a.pdf}
  \includegraphics[width=0.4\textwidth]{figures/example-image-b.pdf}
  \caption{图片示例}
\end{figure}

\printbibliography
\appendix
\chapter{XX代码}
\begin{acknowledgments}
  谢辞应以简短的文字对课题研究与论文撰写过程中曾直接给予帮助的人员
  （例如指导教师、答疑教师及其他人员）表示对自己的谢意，这不仅是一种礼
  貌，也是对他人劳动的尊重，是治学者应当遵循的学术规范。内容限一页。
\end{acknowledgments}

\end{document}
